\documentclass[10pt,a4paper]{amsart}
\usepackage{amsmath,amssymb,amsthm,graphicx}
\usepackage{verbatim}
\newenvironment{reference}{\comment}{\endcomment}
%\newenvironment{reference}{}{}
\newenvironment{slogan}{\comment}{\endcomment}
\newenvironment{history}{\comment}{\endcomment}
% For commutative diagrams we use Xy-pic
\usepackage[all]{xy}

% We use 2cell for 2-commutative diagrams.
\xyoption{2cell}
% Theorem environments.
%
\theoremstyle{plain}
\newtheorem{theorem}[subsection]{Theorem}
\newtheorem{proposition}[subsection]{Proposition}
\newtheorem{lemma}[subsection]{Lemma}

\theoremstyle{definition}
\newtheorem{definition}[subsection]{Definition}
\newtheorem{example}[subsection]{Example}
\newtheorem{exercise}[subsection]{Exercise}
\newtheorem{situation}[subsection]{Situation}

\theoremstyle{remark}
\newtheorem{remark}[subsection]{Remark}
\newtheorem{remarks}[subsection]{Remarks}

\numberwithin{equation}{subsection}

% Macros
%
\def\lim{\mathop{\mathrm{lim}}\nolimits}
\def\colim{\mathop{\mathrm{colim}}\nolimits}
\def\Spec{\mathop{\mathrm{Spec}}}
\def\Hom{\mathop{\mathrm{Hom}}\nolimits}
\def\Ext{\mathop{\mathrm{Ext}}\nolimits}
\def\SheafHom{\mathop{\mathcal{H}\!\mathit{om}}\nolimits}
\def\SheafExt{\mathop{\mathcal{E}\!\mathit{xt}}\nolimits}
\def\Sch{\mathit{Sch}}
\def\Mor{\mathop{\mathrm{Mor}}\nolimits}
\def\Ob{\mathop{\mathrm{Ob}}\nolimits}
\def\Sh{\mathop{\mathit{Sh}}\nolimits}
\def\NL{\mathop{N\!L}\nolimits}
\def\CH{\mathop{\mathrm{CH}}\nolimits}
\def\proetale{{pro\text{-}\acute{e}tale}}
\def\etale{{\acute{e}tale}}
\def\QCoh{\mathit{QCoh}}
\def\Ker{\mathop{\mathrm{Ker}}}
\def\Im{\mathop{\mathrm{Im}}}
\def\Coker{\mathop{\mathrm{Coker}}}
\def\Coim{\mathop{\mathrm{Coim}}}

% Boxtimes
%
\DeclareMathSymbol{\boxtimes}{\mathbin}{AMSa}{"02}

%
% Macros for moduli stacks/spaces
%
\def\QCohstack{\mathcal{QC}\!\mathit{oh}}
\def\Cohstack{\mathcal{C}\!\mathit{oh}}
\def\Spacesstack{\mathcal{S}\!\mathit{paces}}
\def\Quotfunctor{\mathrm{Quot}}
\def\Hilbfunctor{\mathrm{Hilb}}
\def\Curvesstack{\mathcal{C}\!\mathit{urves}}
\def\Polarizedstack{\mathcal{P}\!\mathit{olarized}}
\def\Complexesstack{\mathcal{C}\!\mathit{omplexes}}
% \Pic is the operator that assigns to X its picard group, usage \Pic(X)
% \Picardstack_{X/B} denotes the Picard stack of X over B
% \Picardfunctor_{X/B} denotes the Picard functor of X over B
\def\Pic{\mathop{\mathrm{Pic}}\nolimits}
\def\Picardstack{\mathcal{P}\!\mathit{ic}}
\def\Picardfunctor{\mathrm{Pic}}
\def\Deformationcategory{\mathcal{D}\!\mathit{ef}}
\usepackage{cleveref}
\crefname{originalreference}{original reference}{original references}
\Crefname{originalreference}{Original reference}{Original references}
\makeatletter
\newlabel{categories-section-inertia}{{categories-section-inertia}{source}{}{sourcepointer-1}{}}
\newlabel{categories-section-inertia@cref}{{[originalreference][0][]categories-section-inertia}{[1][0][]source}{}{}{}}
\makeatother
\crefname{originalreference}{original reference}{original references}
\Crefname{originalreference}{Original reference}{Original references}
\makeatletter
\newlabel{stacks-properties-lemma-monomorphism}{{stacks-properties-lemma-monomorphism}{source}{}{sourcepointer-2}{}}
\newlabel{stacks-properties-lemma-monomorphism@cref}{{[originalreference][0][]stacks-properties-lemma-monomorphism}{[1][0][]source}{}{}{}}
\makeatother
\title{The square relating the inertias of a monomorphism of algebraic stacks is cartesian}
\author{The Stacks Project Authors}
\date{stacks-project:stacks-morphisms.tex:local-snapshot-20261012; GNU FDL 1.2 or later}
\begin{document}
\maketitle
\noindent\textit{Editorial note.} All delivered mathematical text is the author's own, copied from the pinned Stacks Project chapter source sources/stacks-algebraic-stacks-morphisms/pinned-original-source.tex; no statement or proof step was rewritten, reordered or removed.\par\medskip
\noindent\textit{Editorial note.} Preamble provenance: the Stacks Project keeps one single shared preamble for every chapter; this card takes it directly from the staged sources/stacks-algebraic-stacks-morphisms/preamble.tex as three verbatim spans (lines 8-12, 14-18 and 154-217), so no derived file was written and nothing was added or redefined. The wrapper class and the base packages are supplied by the assembler.\par\medskip
\noindent\textit{Editorial note.} Bibliography provenance: neither the counted statement nor its proof contains a citation (this lane rescans the whole span before writing the card), so the card carries no bibliography and no entry was reconstructed.\par\medskip
\noindent\textit{Editorial note.} Selection provenance: this card was opened under the file-based control contract (CONTROL.md), with the mechanical pre-card gate run on its statement and proof interval before the card was written.\par\medskip
\noindent\textit{Allowed standard prerequisites (editorial).}\begin{itemize}
\item The characterisation of monomorphisms of algebraic stacks through full faithfulness, and the description of the inertia stack through objects together with their automorphisms. Invoked at the start of the counted proof (source lines 1495-1498) as pointers to the chapter on properties of stacks and to the chapter on categories. Turns an object of the inertia into a pair and lets automorphisms of the source be identified with automorphisms of the image.
\item The description of objects of a fibre product of stacks over a scheme as triples consisting of objects and an isomorphism. Used at the end of the counted proof (source lines 1505-1509). Identifies the objects computed in the proof with objects of the fibre product appearing in the statement.
\end{itemize}\medskip
\ifdefined\hypertarget\hypertarget{sourcepointer-1}{}\fi
\noindent\textit{Source reference.} Original reference categories-section-inertia: sources/stacks-algebraic-stacks-morphisms/pinned-original-source.tex, source lines 1217--1217. Navigation pointer only: the counted text refers to the Stacks item categories-section-inertia, which belongs to another chapter of the project and is not part of this card. The pointer carries the target's own label as its number because this chapter's staged snapshot does not contain the other chapter's tag table; its content is not reproduced here.\par\medskip
\ifdefined\hypertarget\hypertarget{sourcepointer-2}{}\fi
\noindent\textit{Source reference.} Original reference stacks-properties-lemma-monomorphism: sources/stacks-algebraic-stacks-morphisms/pinned-original-source.tex, source lines 1167--1167. Navigation pointer only: the counted text refers to the Stacks item stacks-properties-lemma-monomorphism, which belongs to another chapter of the project and is not part of this card. The pointer carries the target's own label as its number because this chapter's staged snapshot does not contain the other chapter's tag table; its content is not reproduced here.\par\medskip
\begin{lemma}
\label{lemma-monomorphism-cartesian-square-inertia}
Let $f : \mathcal{X} \to \mathcal{Y}$ be a monomorphism of algebraic stacks.
Then the diagram
$$
\xymatrix{
\mathcal{I}_\mathcal{X} \ar[r] \ar[d] &
\mathcal{X} \ar[d] \\
\mathcal{I}_\mathcal{Y} \ar[r] &
\mathcal{Y}
}
$$
is a fibre product square.
\end{lemma}

\begin{proof}
This follows immediately from the fact that $f$ is fully faithful (see
Properties of Stacks, Lemma \ref{stacks-properties-lemma-monomorphism})
and the definition of the inertia in
Categories, Section \ref{categories-section-inertia}.
Namely, an object of $\mathcal{I}_\mathcal{X}$ over a scheme $T$ is
the same thing as a pair $(x, \alpha)$ consisting of an object
$x$ of $\mathcal{X}$ over $T$ and a morphism $\alpha : x \to x$ in
the fibre category of $\mathcal{X}$ over $T$. As $f$ is fully faithful
we see that $\alpha$ is the same thing as a morphism
$\beta : f(x) \to f(x)$ in the fibre category of $\mathcal{Y}$ over $T$.
Hence we can think of objects of $\mathcal{I}_\mathcal{X}$ over $T$
as triples $((y, \beta), x, \gamma)$ where $y$ is an object of
$\mathcal{Y}$ over $T$, $\beta : y \to y$ in $\mathcal{Y}_T$ and
$\gamma : y \to f(x)$ is an isomorphism over $T$, i.e., an object
of $\mathcal{I}_\mathcal{Y} \times_\mathcal{Y} \mathcal{X}$ over $T$.
\end{proof}

\end{document}
